\section{Vhost-Guest Device}\label{sec:Device Types / Vhost-Guest Device}

The virtio vhost-guest device facilitates vhost-user device
emulation through vhost-user protocol exchanges and access to
shared memory. Software-defined networking, storage, and other
I/O appliances can provide services through the device.

This section relies on definitions from the
\hyperref[intro:Vhost-user Protocol]{vhost-user protocol}.
Knowledge of the vhost-user protocol is a prerequisite for
understanding this device.

The \hyperref[intro:Vhost-user Protocol]{vhost-user protocol}
was originally designed for processes on a single system
communicating over UNIX domain sockets. The vhost-guest device
allows the vhost-user backend to communicate with the vhost-user
frontend over the device instead of a UNIX domain socket. This
allows the backend and frontend to run on two separate systems,
such as a virtual machine and a hypervisor. The vhost-user
backend program exchanges vhost-user protocol messages with the
vhost-user frontend through the device.

How the device implementation communicates with the vhost-user
frontend is beyond the scope of this specification. However,
it is expected that most implementations will use an AF_UNIX
socket to relay messages to a vhost-user frontend process
running on the same host.

Existing vhost-user backend programs that communicate over UNIX
domain sockets can support the vhost-guest device without
invasive changes because the pre-existing vhost-user wire
protocol is used. Vhost-user clients do not require modification.
It is also possible to implement a relay that uses a vhost-guest
device and listens on a vhost-user socket. This allows using
unmodified vhost-user servers.

\subsection{Terminology and Typical Operation}\label{sec:Device Types / Vhost-Guest Device / Terminology and Typical Operation}

A system using vhost-guest will generally have multiple
components arranged as follows:

\begin{description}
\item[Backend guest] The guest VM that is providing a vhost-user
    device.

\item[Backend server] A userspace process in the backend guest.
    It transmits and receives vhost-user protocol messages via
    the vhost-guest driver. It has direct access to the memory
    shared by the frontend and may be able to send interrupts
    directly to the frontend. This process does not always have
    administrative privileges in the backend guest.

\item[Vhost-guest driver] A kernel-mode driver in the backend
    guest. It proxies messages between the backend server and the
    vhost-guest device. It also provides an interface that the
    backend server can use to map the frontend guest's memory.

\item[Vhost-guest device] A virtio device implemented in a
    userspace process on the host, typically the VMM running the
    backend guest. It implements a vhost-user backend that
    listens for connections on an AF_UNIX socket. It relays
    messages between the vhost-user frontend and the vhost-guest
    driver. It also exposes the frontend guest's memory to the
    vhost-guest driver.

\item[Vhost-user frontend] A component of the userspace VMM that
    runs the frontend guest. It exposes a virtio device to the
    frontend guest and uses the vhost-user protocol to
    communicate with the vhost-guest device.

\item[Frontend guest] The guest that uses the virtio device
    implemented by the backend server.
\end{description}

Control messages, such as configuration space access and memory
mapping, typically traverse all of the above components. Access
to frontend guest memory typically bypasses the backend server,
vhost-guest driver, vhost-guest device, and vhost-user frontend.
Interrupts typically bypass the vhost-user frontend and
vhost-guest device. They may bypass the vhost-guest driver in
some configurations.

The above component arrangement is not required by this
specification. However, understanding it is necessary to
understand design choices made in this specification.

When this specification refers to the vhost-user device being
implemented, this will always be identified via terms like
``vhost-user'', ``implemented'', or ``frontend.''

Any part of this specification that does not explicitly mention
the implemented device refers to the vhost-guest device itself.

\subsection{Device ID}\label{sec:Device Types / Vhost-Guest Device / Device ID}
  53

\subsection{Feature bits}\label{sec:Device Types / Vhost-Guest Device / Feature bits}

The following feature bits are defined for the vhost-guest device
itself:

\begin{description}

\item[VIRTIO_VHOST_GUEST_F_BIG_ENDIAN (0)]
    If this feature bit is offered, vhost-user protocol messages
    use big-endian byte order. If it is not offered, vhost-user
    protocol messages use little-endian byte order.

    \begin{note}
    Typically, a driver running in a little-endian VM will refuse
    to operate if VIRTIO_VHOST_GUEST_F_BIG_ENDIAN is offered, and
    a driver running in a big-endian VM will refuse to operate if
    VIRTIO_VHOST_GUEST_F_BIG_ENDIAN is not offered. This means
    that the driver must be of the same endianness as the device.
    This is almost always the case in practice.
    \end{note}
\end{description}

Most feature negotiation instead happens via vhost-user's own
mechanism. Details can be found in
\ref{sec:Device Types / Vhost-Guest Device / Restrictions on Vhost-User Protocol Feature Negotiation}.

\subsection{Device Configuration Layout}\label{sec:Device Types / Vhost-Guest Device / Device Configuration Layout}

All fields are always available.

\begin{lstlisting}
struct virtio_vhost_guest_config {
        /* Device implementation limits. */
        le32 max_vhost_queues;
        le32 max_message_size;
        le32 max_migration_message_size;
        le16 max_chained_rx_descriptors;
        le16 max_chained_tx_descriptors;

        /* Which device is this? */
        u8 uuid[16];

        /* Information about the particular device. */
        le16 vhost_user_device_type;
        le16 status;
};
\end{lstlisting}

\begin{description}
\item[\field{status}] contains the status of the vhost-user connection.
    The default value of this field is 0.

    The driver sets VIRTIO_VHOST_GUEST_STATUS_BACKEND_UP to
    indicate readiness for the vhost-user frontend to connect.
    The vhost-user frontend cannot connect until the driver
    sets this bit. Typical implementations will delay either
    bind() or accept() until this point.

    The device sets VIRTIO_VHOST_GUEST_STATUS_FRONTEND_UP when the
    vhost-user frontend connects. After this bit is set, the
    driver can start sending and receiving vhost-user protocol
    messages.

    When the driver clears VIRTIO_VHOST_GUEST_BACKEND_UP while
    the vhost-user frontend is connected, the vhost-user
    frontend is disconnected. This will typically force the
    frontend device into a NEEDS_RESET state.

    When the vhost-user frontend disconnects,
    both VIRTIO_VHOST_GUEST_STATUS_BACKEND_UP and
    VIRTIO_VHOST_GUEST_STATUS_FRONTEND_UP are cleared by the
    device. Communication can be restarted by the driver
    setting VIRTIO_VHOST_GUEST_STATUS_BACKEND_UP again.

    A configuration change notification is sent when the
    device changes the status field, unless the change was caused
    by a write by the driver to the field itself.

\item[\field{max_vhost_queues}] is the maximum number of queues
    that the vhost-user device may have. It is not related to the
    number of queues supported by the vhost-guest device and is
    generally much larger. Devices SHOULD set this field to the
    largest number of virtqueues supported by the vhost-user
    specification, unless resource constraints require a smaller
    value.

\item[\field{max_message_size}] is the maximum size of a
    vhost-user message supported by this device. This field is
    always at least 512.

\item[\field{vhost_user_device_type}] is the expected type of
    the vhost-user device. 0 is special and indicates that any
    type is permissible.

\item[\field{uuid}] is the Universally Unique Identifier (UUID)
    for this device. If the device has no UUID, then this
    field contains the nil UUID (all zeroes). The UUID
    allows vhost-user backend programs to identify a specific
    vhost-guest device among many without relying on bus
    addresses.

\item[\field{max_migration_message_size}] is the maximum size of
    a message allowed on the migration queues (described below).

\end{description}

\devicenormative{\subsubsection}{Device Configuration Layout}{Device Types / Vhost-Guest Device / Device Configuration Layout}

Devices MUST offer VIRTIO_F_VERSION_1 to indicate that the
vhost-guest device is not a legacy device. Devices MUST set
DEVICE_NEEDS_RESET if it is not accepted.

Devices MUST support using reads of any size not exceeding 32
bits to access the configuration space, provided they are
naturally aligned. Such reads MUST return an atomic snapshot of
the configuration space. Devices MUST support using 32-bit
writes to write to both \field{status} and
\field{vhost_user_device_type} in a single atomic operation.

Devices MUST support 1-byte reads and writes to the bytes in
\field{flags}.

\field{max_chained_tx_descriptors} and
\field{max_chained_rx_descriptors} MUST both be at least 1.

Devices MUST NOT change \field{uuid} except when the device is
reset.

\drivernormative{\subsubsection}{Device Configuration Layout}{Device Types / Vhost-Guest Device / Device Configuration Layout}

The driver MUST refuse to use the device if VIRTIO_F_VERSION_1 is
not offered. The driver MUST accept VIRTIO_F_VERSION_1 if it is
offered.

The driver MUST NOT write to device configuration fields other
than \field{status} and \field{vhost_user_device_type}.

The driver MUST NOT chain more than
\field{max_chained_tx_descriptors} device-readable descriptors.
The driver MUST NOT chain more than
\field{max_chained_rx_descriptors} device-writeable descriptors.
If either of these fields are 0, the driver MUST refuse to
operate the device.

\subsection{Virtqueues}\label{sec:Device Types / Vhost-Guest Device / Virtqueues}

The following virtqueues are always available:

\begin{description}
\item[0] front2back_requests (device-to-driver vhost-user requests)
\item[1] front2back_replies (driver-to-device vhost-user replies)
\item[2] back2front_replies (device-to-driver vhost-user replies)
\item[3] back2front_requests (driver-to-device vhost-user requests)
\end{description}

If the implemented device has negotiated the
VHOST_USER_PROTOCOL_F_DEVICE_STATE feature, the following
virtqueues are available:

\begin{description}
\item[4] migration_send (migration state sending)
\item[5] migration_send (migration state receiving)
\end{description}

\drivernormative{\subsubsection}{Virtqueues}{Device Types / Vhost-Guest Device / Virtqueues}

The driver MUST NOT access the \field{migration_save} virtqueue or
the \field{migration_load} virtqueue unless the implemented
device has negotiated the VHOST_USER_PROTOCOL_F_DEVICE_STATE
feature.

The driver MUST NOT access the \field{migration_save} virtqueue
unless a VHOST_USER_SET_DEVICE_STATE_FD message with direction
set to 0 (Save) has been received.

The driver MUST NOT access the \field{migration_load} virtqueue
unless a VHOST_USER_SET_DEVICE_STATE_FD message with direction
set to 1 (Load) has been received.

\subsection{Device Initialization}\label{sec:Device Types / Vhost-Guest Device / Device Initialization}

Initialization is partly described in
\ref{sec:Device Types / Vhost-Guest Device / Device Configuration Layout}
and partly in the normative statements below.

\drivernormative{\subsubsection}{Device Initialization}{Device Types / Vhost-Guest Device / Device Initialization}

The driver MUST fail feature negotiation if VIRTIO_F_VERSION_1 is
not offered. The driver MUST accept VIRTIO_F_VERSION_1 if it is
offered.

If the VIRTIO_VHOST_GUEST_F_BIG_ENDIAN feature is offered by the
device, the driver MUST accept it or refuse to finish feature
negotiation.

If the VIRTIO_VHOST_GUEST_F_BIG_ENDIAN feature is not offered by
the device but is required by the driver, the driver MUST refuse
to finish feature negotiation.

\begin{note}
Support for endian swapping in the device was considered but not
included. The backend needs to be aware of each vhost-user
protocol message, whereas the vhost-guest device does not.
Therefore, byte-swapping in the backend is significantly simpler.
The vast majority of use-cases don't require byte-swapping at
all. Furthermore, supporting byte-swapping in the device would
make feature negotiation significantly more complex.
\end{note}

The driver MUST read from \field{vhost_user_device_type} before
writing to it or to \field{status}. If the read returns 0, the
driver MUST write a nonzero value to
\field{vhost_user_device_type} before writing to \field{status},
or use a 32-bit write to write both at once. If the read returns
nonzero, the driver MUST NOT write to
\field{vhost_user_device_type}. The driver MUST do this before
setting VIRTIO_VHOST_GUEST_BACKEND_UP.

If the driver sets VIRTIO_VHOST_GUEST_BACKEND_UP, the driver MUST
implement a device that conforms to the specification for a
device of the type it read from vhost_user_device_type (if that
is nonzero), or of the type that it wrote to
vhost_user_device_type (if reading it returned zero). If the
driver cannot do this (perhaps because the device type is one it
does not implement), it MUST refuse to operate the device.

The driver SHOULD check the \field{max_vhost_queues}
configuration field to determine how many queues the
vhost-user backend will be able to support.

If a driver supports multiple vhost-guest the device and needs
to differentiate them, it SHOULD use the \field{uuid}
configuration field to allow vhost-user backend programs to
identify a specific device among many.

The driver SHOULD place at least one buffer in
front2back_requests before setting the
VIRTIO_VHOST_GUEST_BACKEND_UP bit in the \field{status}
configuration field.

The driver MUST handle virtqueue messages on front2back_requests
notifications that occur before the configuration change
notification. It is possible that a vhost-user protocol
message from the vhost-user frontend arrives before the
driver has seen the configuration change notification for the
VIRTIO_VHOST_GUEST_STATUS_FRONTEND_UP \field{status} change.

\subsection{Device Operation}\label{sec:Device Types / Vhost-Guest Device / Device Operation}

Device operation consists of operating request queues, response
queues, and migration queues.

\subsubsection{Device Operation: Request And Response Queues}\label{sec:Device Types / Vhost-Guest Device / Device Operation / Device Operation: Request And Response Queues}

The driver receives vhost-user requests from the vhost-user
frontend on front2back_requests. The driver sends responses
to the vhost-user frontend on back2front_replies.

The driver sends vhost-user requests to the vhost-user frontend
on back2front_requests. The driver receives responses from
the vhost-user frontend on front2back_replies.

All descriptors on front2back_requests and front2back_replies are
write-only to the device. All descriptors on back2front_requests
and back2front_replies are read-only to the device.

Unless stated elsewhere in this specification, messages from the
frontend to backend and backend to frontend are delivered without
modification.

Within a virtqueue, all messages are guaranteed to be delivered
in-order. There is no ordering between virtqueues.

Each buffer is a vhost-user protocol message as defined by
the \hyperref[intro:Vhost-user Protocol]{vhost-user protocol}.

If the VIRTIO_VHOST_GUEST_F_BIG_ENDIAN feature is negotiated,
messages are in big-endian order. Otherwise, they are in
little-endian order.

\drivernormative{\paragraph}{Device Operation: Request And Response Queues}{Device Types / Vhost-Guest Device / Device Operation / Device Operation: Request And Response Queues}

Descriptors placed on \field{back2front_requests},
\field{back2front_replies}, and \field{migration_save} MUST be
device-readable.

Descriptors placed on \field{front2back_requests},
\field{front2back_replies}, and \field{migration_load} MUST be
device-writeable.

Descriptors placed on front2back_requests and front2back_replies
MUST have a length not less than the smaller of:

\begin{enumerate}
\item The contents of the \field{max_message_size}
      configuration field.
\item 4KiB.
\end{enumerate}

Descriptors placed on back2front_requests and back2front_replies
MUST have a length that does not exceed \field{max_message_size}.

The driver MUST ensure that each descriptor on back2front_requests
and back2front_replies points to a valid and complete vhost-user
protocol message that is appropriate for the queue that it is on.
Padding at the end of a descriptor is not allowed.

\devicenormative{\paragraph}{Device Operation: Request And Response Queues}{Device Types / Vhost-Guest Device / Device Operation / Device Operation: Request And Response Queues}

Devices MUST ensure that the header of each message in
\field{front2back_requests} and \field{front2back_replies}
conforms to the
\hyperref[intro:Vhost-user Protocol]{vhost-user protocol}.
Violations MUST cause the device to set DEVICE_NEEDS_RESET.

Devices MAY perform such checking on the body of the message. If
they detect that the
\hyperref[intro:Vhost-user Protocol]{vhost-user protocol} has
been violated, they MAY set DEVICE_NEEDS_RESET.

If a vhost-user protocol message does not fit in the buffer
offered by the driver, the device MUST set DEVICE_NEEDS_RESET.

\subsubsection{Device Operation: Migration Queues}\label{sec:Device Types / Vhost-Guest Device / Device Operation / Device Operation: Migration Queues}

The driver receives incoming migration data on migration_receive.
The driver sends migration data on migration_send.

All descriptors on migration_receive are write-only to the
device. All descriptors on migration_send are read-only to the
device.

Descriptors written to the migration_send queue are written to
the migration pipe without modification. The device closes the
migration pipe when it gets an empty descriptor.

Data read from the inbound migration data pipe is written to
descriptors on the migration_receive queue. EOF or error is
indicated by writing zero bytes to a descriptor.

\drivernormative{\paragraph}{Device Operation: Migration Queues}{Device Types / Vhost-Guest Device / Device Operation / Device Operation: Migration Queues}

The driver MUST limit messages on \field{migration_send} to the
value in the configuration field
\field{max_migration_message_size}.

Devices silently truncate the migration stream if an error
occurs. If the driver did not use a prefix-free encoding for
the migration stream, it would not be able to detect errors.
Therefore, the driver MUST use a prefix-free encoding for
migration streams.

\subsubsection{Device Operation: File Descriptor Passing}\label{sec:Device Types / Vhost-Guest Device / Device Operation / Device Operation: File Descriptor Passing}

File descriptor passing is handled differently by the vhost-user
device backend. When a message is received that carries one
or more file descriptors according to the vhost-user protocol,
additional device resources become available to the driver, as
described in \ref{sec:Device Types / Vhost-Guest Device / Additional Resources / Availability of Additional Resources}.

\subsection{Additional Resources}\label{sec:Device Types / Vhost-Guest Device / Additional Resources}

The vhost-guest device uses the following additional features of
a virtio device beyond virtqueues and configuration space:

\subsubsection{Device Auxiliary Notifications}\label{sec:Device Types / Vhost-Guest Device / Additional Resources / Device Auxiliary Notifications}

Device auxiliary notifications are laid out as follows:

\begin{description}
\item[0] Vring call for vhost-user queue 0
\item[\ldots]
\item[N] Vring err for vhost-user queue 0
\item[\ldots]
\item[2N] Log
\end{description}

\subsubsection{Driver Auxiliary Notifications}\label{sec:Device Types / Vhost-Guest Device / Additional Resources / Driver Auxiliary Notifications}

Driver auxiliary notifications are laid out as follows:

\begin{description}
\item[0] Vring kick for vhost-user queue 0
\item[\ldots]
\end{description}

\subsubsection{Shared Memory}\label{sec:Device Types / Vhost-Guest Device / Additional Resources / Shared Memory}

Shared memory offsets are chosen by the device. The device
promises that they do not overlap. The device replaces the user
address field in a vhost-user protocol message with the offset in
the shared memory region. If a message has a guest address but no
user address, the guest address is replaced instead.

The driver MAY detect overlap. If it does, it SHOULD set FAILED.

\subsubsection{Availability of Additional Resources}\label{sec:Device Types / Vhost-Guest Device / Additional Resources / Availability of Additional Resources}

The following vhost-user protocol messages convey access to
additional device resources:

\begin{description}
\item[VHOST_USER_SET_MEM_TABLE]
    Contents of vhost memory regions are available to the
    driver in shared memory. The user addresses in this message
    have been replaced with offsets in the shared memory region.
\item[VHOST_USER_ADD_MEM_REG]
    A new vhost memory region is available to the driver in
    shared memory. The user address in this message has been
    replaced with an offset in the shared memory region.
\item[VHOST_USER_REM_MEM_REG]
    The frontend wishes to remove the vhost memory region
    described by this message. The memory will remain available
    until the driver replies to this message.
\item[VHOST_USER_SET_LOG_BASE]
    Contents of the log are available to the driver in shared
    memory. The offset has been replaced with the offset in the
    shared memory region.
\item[VHOST_USER_SET_LOG_FD]
    The log device auxiliary notification is available to the
    driver. Using this notification before this message is
    received has no effect.
\item[VHOST_USER_GET_INFLIGHT_FD]
    The inflight memory region is available to the driver. The
    mmap offset has been replaced with the offset in the ring
    buffer.
\item[VHOST_USER_SET_INFLIGHT_FD]
    The inflight memory region is available to the driver. The
    mmap offset has been replaced with the offset in the ring
    buffer.
\item[VHOST_USER_SET_VRING_KICK]
    The vring kick driver auxiliary notification for this queue
    is available to the driver. If this notification is enabled,
    it may trigger before the driver has processed this message.
\item[VHOST_USER_SET_VRING_CALL]
    The vring call device auxiliary notification for this queue
    is available to the driver. Using this notification before
    before this message is received has no effect.
\item[VHOST_USER_SET_VRING_ERR]
    The vring err device auxiliary notification for this queue
    is available to the driver. Using this notification before
    before this message is received has no effect.
\item[VHOST_USER_SET_BACKEND_REQ_FD]
    The driver may send vhost-user protocol backend messages on
    back2front_requests. Putting buffers on back2front_requests
    before this message is received is not allowed.
\item[VHOST_USER_SET_DEVICE_STATE_FD]
    The driver may send its migration state on migration_send or
    receive them on migration_receive. This state consists of an
    opaque stream of bytes. The first buffer may be placed on
    migration_receive, and its corresponding notification
    received by the driver, before this message arrives.
\end{description}

\drivernormative{\paragraph}{Availability of Additional Resources}{Device Types / Vhost-Guest Device / Additional Resources / Availability of Additional Resources}

The driver MUST NOT place a buffer on back2front_requests until
and unless VHOST_USER_SET_BACKEND_REQ_FD is received.

The driver MUST NOT place a buffer on migration_send until and
unless VHOST_USER_SET_DEVICE_STATE_FD is received with a
transfer direction of ``load''.

\devicenormative{\paragraph}{Availability of Additional Resources}{Device Types / Vhost-Guest Device / Additional Resources / Availability of Additional Resources}

Devices MUST ensure all memory regions have non-overlapping
$\left[user address, user address + size\right)$.

Devices MUST ensure all memory regions have non-overlapping
$\left[guest address, guest address + size\right)$.

\subsubsection{Notifications}\label{sec:Device Types / Vhost-Guest Device / Additional Resources / Notifications}

When the device has made a buffer available for the driver
to process, it notifies the driver via a driver auxiliary
notification. Therefore, only transports supporting driver
auxiliary notifications are supported by this device.

When the driver has made a buffer available for the device to
process, it notifies the device via a device auxiliary
notification. Therefore, only transports supporting device
auxiliary notifications are supported by this device.

\subsection{Restrictions on Vhost-User Protocol Feature Negotiation}\label{sec:Device Types / Vhost-Guest Device / Restrictions on Vhost-User Protocol Feature Negotiation}

To ensure forward compatibility with future versions of the
vhost-guest device, the set of vhost-user protocol features
that can be negotiated is restricted.

While many vhost-user protocol messages can be forwarded
unchanged, not all can. For instance, vhost-user messages that
carry file descriptors, guest addresses, or user addresses
require special handling by the device. Devices may choose to not
implement all of these messages. Furthermore, a device that does
not know about a protocol feature does not know if it requires
any special handling. Devices ensure that only supported
protocol features are used by clearing protocol bits in
VHOST_USER_GET_PROTOCOL_FEATURES replies and by validating
VHOST_USER_SET_PROTOCOL_FEATURES against the protocol features
supported by the device

\devicenormative{\subsubsection}{Restrictions on Vhost-User Protocol Feature Negotiation}{Device Types / Vhost-Guest Device / Restrictions on Vhost-User Protocol Feature Negotiation}

A device MUST clear bits in VHOST_USER_GET_PROTOCOL_FEATURES
replies that correspond to features it does not support. It MAY
clear bits in VHOST_USER_GET_PROTOCOL_FEATURES replies for other
reasons. A device MUST NOT modify the contents of a
VHOST_USER_GET_PROTOCOL_FEATURES reply in any other way.

If a device receives a VHOST_USER_SET_PROTOCOL_FEATURE message
with a feature bit set that corresponds to a feature it does not
support, it MUST disconnect the frontend. If a device receives a
VHOST_USER_SET_PROTOCOL_FEATURE message with a feature bit set
that corresponds to a feature it does not permit for other
reasons, it MAY disconnect the frontend. Devices MUST NOT modify
VHOST_USER_SET_PROTOCOL_FEATURE messages.

\subsection{Frontend Disconnection}\label{sec:Device Types / Vhost-Guest Device / Frontend Disconnection}

When the device disconnects from the frontend, the device sets
DEVICE_NEEDS_RESET. Shared memory regions remain intact until the
driver resets the device.

\drivernormative{\subsubsection}{Frontend Disconnection}{Device Types / Vhost-Guest Device / Frontend Disconnection}

If the driver wants to keep using the device, it MUST stop
accessing the shared memory regions, then reset the device.

\devicenormative{\subsubsection}{Frontend Disconnection}{Device Types / Vhost-Guest Device / Frontend Disconnection}

If a device wants the frontend to continue to provide the same
service, it SHOULD NOT change \field{uuid} after being reset.

If a device wants the frontend to provide a different service,
it MUST change \field{uuid}.

\begin{note}

This implies that if the device does not have a UUID, the
service it provides cannot be changed except by hot-unplugging
it.

\end{note}
